\chapter{Conclusiones y Recomendaciones}
\label{sec.tp2:conclusiones}

Este capítulo compila la evaluación conclusiva de la instalación de iluminación relevada en base a las mediciones
registradas y al análisis comparativo con el marco legal.

\section{Conclusión del Relevamiento}
\label{sec.tp2:conclusion_evaluacion}

Luego de procesar y evaluar las mediciones realizadas en el local, se determinan las siguientes conclusiones técnicas:
\begin{itemize}
    \item \textbf{Sector Oficina:} La iluminancia media general es satisfactoria ($658{,}89\lx \ge 300\lx$), pero \textbf{no cumple con el criterio de uniformidad legal} debido al fuerte contraste provocado por la radiación solar directa que ingresa por las ventanas. Las zonas más alejadas de las aberturas dependen únicamente del plafón central de $18\W$ y registran niveles insuficientes ($188\lx$), por lo que existe riesgo de fatiga visual y deslumbramiento.
    \item \textbf{Sector Depósito:} El sector \textbf{cumple plenamente} con el marco legal (Dec. N° 351/79), tanto en nivel medio ($197{,}78\lx \ge 100\lx$) como en uniformidad ($140\lx \ge 98{,}89\lx$). La iluminación artificial general es adecuada y homogénea.
    \item \textbf{Escritorio Largo (Puesto Localizado):} El puesto \textbf{no cumple} con las exigencias legales para dibujo y diseño técnico ($345{,}67\lx < 500\lx$). La tira LED de $5050$ actual es insuficiente a la altura de montaje de $62\cm$.
    \item \textbf{Escritorio Chico (Puesto Localizado):} El puesto \textbf{cumple satisfactoriamente} con el nivel exigido de $500\lx$, alcanzando una iluminancia de $710\lx$ gracias al correcto funcionamiento y direccionamiento de la luminaria de mesa regulable de $700\lm$.
\end{itemize}

\section{Medidas de Recomendación y Mejora}
\label{sec.tp2:recomendaciones}

Con el fin de subsanar los desvíos detectados y optimizar las condiciones ergonómicas y de seguridad del local, se recomiendan las siguientes acciones correctivas y preventivas:

\begin{enumerate}
    \item \textbf{Readecuación de la Iluminación General de la Oficina:} Se aconseja instalar un segundo plafón LED de $18\W$ en el cielorraso de la oficina, disponiéndolos simétricamente. Esto reforzará los niveles de iluminancia artificial en las zonas alejadas de la ventana (puntos $E_2, E_3, E_4$), elevándolos sobre los $300\lx$ y garantizando una óptima relación de uniformidad ($E_{min} \ge E_{media}/2$).
    \item \textbf{Control de la Luz Natural Directa:} Implementar cortinas translúcidas, persianas regulables o vinilos difusores en las ventanas (en especial la ventana de la derecha). Esto atenuará los picos excesivos de radiación solar directa ($1851\lx$) que causan deslumbramiento y elevados contrastes térmicos y lumínicos.
    \item \textbf{Reemplazo de la Tira LED en el Escritorio Largo:} Sustituir la tira LED $5050$ actual por una tira LED de mayor flujo luminoso y densidad (por ejemplo, una tira de tecnología SMD 2835 o de doble hilera con una densidad de $120\text{ leds}/\m$ o superior). Esto aumentará la iluminancia sobre el plano de trabajo por encima de los $500\lx$ requeridos para actividades visuales de precisión.
    \item \textbf{Programa de Limpieza Periódica:} Establecer un cronograma semestral de limpieza de difusores y lámparas para evitar la pérdida de rendimiento por sedimentación de polvo (depreciación del flujo luminoso).
    \item \textbf{Monitoreo Anual Obligatorio:} Programar nuevas mediciones preventivas con frecuencia anual, de acuerdo con las directrices de la Resolución SRT N° 84/2012, para verificar la constancia de las condiciones lumínicas del ambiente de trabajo.
\end{enumerate}
